\documentclass[12pt,a4paper]{article}
\usepackage[utf8]{inputenc}   % Encodage des caractères en UTF-8
\usepackage[T1]{fontenc}      % Encodage de la police
\usepackage{amsmath,amssymb}  % Symboles mathématiques
\usepackage{geometry}
\usepackage{listings}
\usepackage{color}
\usepackage{tcolorbox}

\thispagestyle{empty}

\geometry{left=1.5cm, right=1.5cm, top=2cm, bottom=2cm}

\newcommand{\listingsttfamily}{\fontfamily{IBMPlexMono-TLF}\small}

\definecolor{mygreen}{rgb}{0,0.6,0}
\definecolor{mygray}{rgb}{0.5,0.5,0.5}
\definecolor{mymauve}{rgb}{0.58,0,0.82}
\definecolor{mywhite}{rgb}{0.9,0.9,0.9}

\lstset{ %
  backgroundcolor=\color{mywhite},   % choose the background color
  %basicstyle=\footnotesize,        % size of fonts used for the code
  basicstyle=\listingsttfamily,
  breaklines=true,                 % automatic line breaking only at whitespace
  captionpos=b,                    % sets the caption-position to bottom
  commentstyle=\color{mygreen},    % comment style
  escapeinside={\%*}{*)},          % if you want to add LaTeX within your code
  keywordstyle=\color{blue},       % keyword style
  stringstyle=\color{mymauve},     % string literal style
}

\begin{document}

\begin{center}
  \Large
  Contrôle - Chiffrement affine en Python.
\end{center}

\vfill

\begin{tcolorbox}
  
\textbf{Contexte :}  
Nous souhaitons chiffrer un message en utilisant un chiffrement affine, où pour chaque caractère
(on le représente par un entier \texttt{x}), on calcule :

\[
y \equiv a \times x + b \pmod{m}.
\]

Le déchiffrement utilise l'inverse de \texttt{a} modulo \texttt{m}, noté \texttt{a\(^{-1}\)}, de sorte que :

\[
x \equiv a^{-1} \times (y - b) \pmod{m}.
\]
\end{tcolorbox}

%Les questions suivantes concernent :
%
%\begin{itemize}
%  \item La modification du message à chiffrer,
%  \item L’adaptation de la plage de caractères (en excluant ou non certains codes ASCII),
%  \item Les formules de chiffrement et de déchiffrement,
%  \item Les tests de validité (inversibilité, portée des caractères, etc.).
%\end{itemize}

\vfill

\begin{enumerate}

  \item \textbf{Changer le message à chiffrer.}\\
  Dans le code Python fourni, expliquez comment remplacer le message par un autre.  
  Quels sont les endroits (variables, lignes) qui doivent être modifiés afin de définir un nouveau texte à chiffrer ?

  % \item \textbf{Conversion d’une liste de codes ASCII en chaîne de caractères.}\\
  % On dispose d’une liste de valeurs entières issues d’un chiffrement, par exemple \verb=[100, 107, 98, ...]=.  
  % Comment la convertir en une chaîne de caractères, par exemple \verb="dkbf..."= ?  
  % (Référence : fonctions \texttt{chr}, \texttt{join}, et boucles \texttt{for}.)
  % % TODO indiquer qu'il peuvent seulement expliquer et non coder

  \item \textbf{Test d'inversibilité de \texttt{a} modulo \texttt{m}.}\\
  On souhaite vérifier que \texttt{a} est bien inversible modulo \texttt{m}.  
  Proposez un code Python minimal utilisant la fonction \texttt{find\_inverse(a, m)}.  
  Affichez \texttt{"OK"} si \texttt{a} est inversible, ou \texttt{"BAD"} sinon.
  % TODO: donner un exemple de if else.

  \item \textbf{Vérification d'une liste chiffrée.}\\
  On reçoit une liste \verb|liste_chiffree| supposée contenir des valeurs comprises entre 0 et \texttt{m-1}.
  \begin{enumerate}
    \item Écrivez un code pour :
      \begin{itemize}
        \item Parcourir cette liste,
        \item Afficher \texttt{"OK"} si l'élément courant est bien dans l'intervalle \texttt{[0 .. m-1]},
        \item Afficher \texttt{"BAD"} sinon.
      \end{itemize}
    \item On souhaite simplifier l'affichage et afficher \texttt{"OK"} si tout les éléments de la liste sont dans l'intervalle \texttt{[0 .. m-1]} et \texttt{"BAD"} sinon. Expliquez une manière de procéder.
  \end{enumerate}


  \item \textbf{Se restreindre au codes ascii affichable.}\\
  On choisi de travailler sur la plage ASCII de 32 (inclus) à 255 (inclus).
  Pour cela on se ramène à des nombres entre 0 et 223 et on travaille modulo $m'=224$.
  \begin{enumerate}
    \item Donner un nouveau $a$ inversible modulo $m'$.
    \item On va appliquer le chiffrement affine modulo $m'$.
    Ainsi lorsqu'on a le code ascii d'un caractère à chiffrer on lui soustrait 32.
    Après le chiffrement on ajoute 32 pour obtenir un code affichable.
    \begin{enumerate}
        \item Expliquez clairement pourquoi on \emph{décale de 32} lors du chiffrement, puis on \emph{réajoute 32} pour l'affichage.
        \item Écrivez la formule de chiffrement. Puis écrire le code python réalisant ce calcul.
        % :
        %\[
        %  y \;=\; 32 + \bigl(a \times (c - 32) + b \bigr)\bmod m',
        %\]
        %où \texttt{c} est le code ASCII du caractère à chiffrer (entre 32 et 255).
        \item Écrivez la formule de déchiffrement.% (en supposant que \texttt{a} est inversible modulo \texttt{m'} et que \texttt{a\(^{-1}\)} est son inverse).
        %\item Indiquez quelles lignes du code Python doivent être modifiées pour utiliser ce nouveau module \texttt{m'} et réaliser correctement les opérations de décalage.
    \end{enumerate}
  \end{enumerate}
  % TODO Joindre code ascii
  \vfill
  \begin{tcolorbox}
    Pour information voici la manière d'écrire des instructions conditionnelles en python :
  \begin{lstlisting}[language=Python]
    if(a % 2 == 0): # le contenu de la variable "a" est-il pair ?
      print("OK") # on affiche "OK" si c'est le cas
    else:
      print("BAD") # on affiche "BAD" sinon
  \end{lstlisting}
  \end{tcolorbox}

  

  % \item \textbf{Restriction au ``ASCII de base'' (32 à 126).}\\
  % On souhaite maintenant se limiter aux caractères dont le code ASCII est compris entre 32 (espace) et 126 (tilde \textasciitilde), soit 95 caractères.
  % \begin{enumerate}
  %   \item Montrez que la taille du nouvel ensemble de caractères est \texttt{m'} = 95.
  %   \item Proposez un entier \texttt{a} qui soit inversible modulo 95. (Autrement dit, $\gcd(a,95)=1$.)
  %   \item Écrivez la formule de chiffrement :
  %   \[
  %     y \;=\; 32 + \bigl(a \times (c - 32) + b \bigr)\bmod 95,
  %   \]
  %   et la formule de déchiffrement correspondante (en utilisant l’inverse de \texttt{a} modulo 95).
  %   \item Expliquez une nouvelle fois quelles modifications (décalages, calculs \texttt{\% 95}, etc.) sont nécessaires dans le code Python.
  % \end{enumerate}

\end{enumerate}

\end{document}
